\section{Reflection}

\subsection{The bottleneck moved}

Thirteen experiments in five working days. At no point was implementation the limiting resource.

What limited the work was \textbf{trust} --- and specifically, one human reading. Every pull request, every ledger entry, every claimed number. The apparatus in \S4 exists because implementation stopped being scarce and verification did not, and a practice built for the old scarcity will optimise the wrong thing.

That inversion has consequences beyond a single project. Research training is organised around production: how to design an experiment, implement it, analyse it, write it up. The implicit model is that thinking is fast and building is slow, so rigour can be afforded because you are waiting for the build anyway. Reverse it --- building becomes fast, checking stays slow --- and rigour stops being something you have time for and becomes the thing you spend your time on.

This is not a lament. It is arguably a better allocation. But it is not the one anyone was trained for, and a researcher who keeps optimising production will produce more than they can vouch for, faster, without ever noticing the moment it happened.

\subsection{Confident and wrong has no analogue in the tools we inherited}

Version control, testing, code review, continuous integration: every one of them was built for artifacts that \textbf{break loudly}. A test fails. A type does not check. A merge conflicts. The entire apparatus of software engineering is a set of amplifiers for a signal that already exists.

A wrong belief emits no signal. It is internally consistent by construction --- that is what makes it believable --- and it passes every check because it is not a defect in an artifact but a defect in what someone concluded from one. The green badge and the wrong claim are perfectly compatible, and in this project they coexisted for ten experiments.

We have almost no tools for this. Pre-registration is one, borrowed from clinical trials and psychology, and it is telling that the borrowing was necessary --- computational research had not needed it while the cost of running an experiment was itself a check on how many you could run and how loosely you could interpret them.

The practices in \S4 are, I think, the beginning of an answer rather than an answer. Non-ratification, provenance on every carried figure, a falsifier for the specification and not just for the hypothesis: these are attempts to make a wrong belief \textit{break loudly}, by manufacturing the signal that does not naturally exist. Whether they are the right attempts is exactly what a second project would find out.

\subsection{The unfashionable deliverable}

The most reusable thing this project produced is not the code, and not its one positive result. It is a ledger of thirteen foreclosed possibilities, each with the evidence that closed it, the conditions under which it holds, and the cost of reopening it.

Someone building a system of this kind now knows several things not to try, and why. That is a different sort of result from a benchmark number, and I would argue a more durable one: a benchmark is superseded by the next method, while a well-founded elimination stays eliminated until someone changes the assumptions it rested on --- which the ledger states, so they can check.

It is also almost unpublishable. There is no venue for \textit{here are thirteen things that do not work and the reasons}, no citation economy that rewards it, and no obvious way to present it at a conference. That is a fact about publishing rather than about the knowledge, and I want to state it plainly rather than pretend the incentives are aligned: \textbf{the better this method works, the more of its output takes a form the field does not reward.}

Which raises a question I cannot answer from one project. If AI collaboration makes elimination cheap --- and it does; most of these thirteen cost an afternoon each --- then a great deal of the total knowledge produced will be of a kind we currently have no way to share. Either the practice changes to accommodate it, or the knowledge stays in repositories that nobody reads.

\subsection{That this report exists at all}

These notes were produced under the method they describe, by the collaboration they analyse, including their own claims-verification pass. The recursion is not decoration. It is either the strongest available evidence for the method or a closed loop, and I owe the reader a position on which.

The honest answer is: \textbf{both, in different parts, and the boundary is checkable.}

The parts that are evidence are the ones where the method's mechanisms fired against this document. The first claims pass ran on the technical report and moved eleven of its eighteen substantive claims. The repeated-figure check fired on the template's own front page. A specification about this work was rejected on a false premise about which surface held which file. Those are not testimonials; they are recorded events with commits attached, and a reader can go and look.

The part that is a closed loop is the argument. \S6's finding --- that a rigorous project will not use its own informal lane --- is my reading of my own behaviour, produced in conversation with the collaborator whose behaviour is also under examination. No mechanism checked it. It survived because it explained something that had happened twice and because I recognised it, which is exactly the standard \S4.1 says is insufficient.

So: trust the events, weigh the arguments. That is the same instruction the method gives about any of its own claims, and applying it here is not modesty but consistency.

\subsection{What I would want checked}

Three things a second project would settle that this one could not.

\textbf{Whether the failures are the method's or mine.} One researcher. Every failure in \S6 is consistent with \textit{this is what happens to a rigorous project} and equally consistent with \textit{this is what happens when this particular person is given a rigorous framework}. I lean toward the first because informal exploration went undone in two consecutive projects --- in the first because there was no lane, in the second although there was --- but two is not many.

\textbf{Whether the mechanisms transfer or were shaped by the domain.} Pre-registration with a falsifier works cleanly when an experiment produces a number that can be compared against a threshold. Whether it works for research whose outputs are qualitative, or exploratory, or where the interesting result is a construction rather than a measurement, I do not know. I would guess the falsifier survives and the threshold does not.

\textbf{Whether the record is readable by anyone but its author.} Nine thousand lines of ledger, written by one collaborator, read by one human, in five working days. A colleague opening it cold is the test that has not been run, and until it is, \textit{the repository is the memory} is a claim about a memory of one.

\subsection{How did it feel}

As stated at the beginning, this is a document written in collaboration between a researcher and an AI entity. The curious reader may be challenged to imagine how this collaboration actually took place---and perhaps even to identify who wrote which parts.

This subsection, starting with its very title, is the only one for which our readership can be certain that it was written by a human---because a machine does not have feelings, right?

So, let us rephrase the question that cries out for an answer: \textbf{``How did it feel to partner with an AI?''} (Which, by the way, I prefer to call an \textit{entity}, not an \textit{agent}.)

Unfortunately, it is indescribable. Like jazz, you have to experience it to know. It is certainly different, and it feels new. It is here to stay, but each one of us, as humans, must find the answer in our own inner way.

On the other hand, for endeavors such as the one described in this document, we can say that the collaboration is a kind of \textit{contract} whose terms the two parties must learn to negotiate in order for things to work out. It is also a moving target, because everything is changing at a rapid pace and requires us to learn as we go.

But since, as a researcher, I am naturally drawn to discovery as well as to the unknown, I can assure you that \textbf{``It Feels Good!''}

\subsection{What I would tell a student starting now}

For a scientist who is also a mentor and educator, this is perhaps the question of a lifetime. It has two inseparable sides: one seen from the perspective of the teacher, and the other from that of the student. Together, they reveal an intrinsic interdependence of purpose, learning, and aspiration.

The present moment is changing not only how we learn, but also what is worth learning. The emphasis is gradually moving away from mechanical procedures and the memorization of rules---the kinds of tasks that machines can now perform remarkably well and with increasing ease. What becomes more important, then, are basic principles, fundamentals, genuine understanding, intuition, judgment, and creativity. And, of course, talent helps.

In that spirit, the advice I would give is simple: ``Discover the things you love and the things you are good at. Find where they meet, and pursue that path.''

Perhaps education, in this new context, is less about preparing someone to follow a predefined path and more about helping them acquire the confidence and understanding needed to discover one of their own.

\subsection{Is this a good way to spend a career?}

To answer this question, we may first need to reconsider what we mean by a \textit{career}. In this new context, perhaps the traditional idea of a career---a more or less predefined trajectory through a profession---will no longer be the most useful one.

A career may increasingly become something that is continuously invented: a sequence of questions, discoveries, collaborations, changes of direction, and new beginnings.

So, to close briefly, I would return to the advice of the previous subsection: discover what you care about, discover what you can contribute, and invent your own career around the intersection of the two.

If you can do that, perhaps the question of whether it is a good way to spend a career becomes less important.

It may simply become a good way to spend a lifetime.
